% Fragment -- inclus de main.tex prin \input
\clearpage
\setpartlabel{themeSeo}{\faShareSquare}{Faza 5: SEO + Open Graph}
\setcounter{section}{0}
\phantomsection
\addcontentsline{toc}{part}{\protect\textbf{Faza 5 \textendash{} SEO + Open Graph}}
\handout{Faza 5 \textendash{} SEO + Open Graph}{Metadate pentru motoarele de cautare si retelele sociale.}

Capitolul prezinta tag-urile meta utilizate de motoarele de cautare (in principal Google) si de platformele sociale (Facebook, X, LinkedIn, WhatsApp, Discord) pentru afisarea corecta a paginilor distribuite. Acoperirea include SEO clasic, protocolul Open Graph (OGP), datele structurate Schema.org, si URL-urile canonice.

\bigskip

\begin{outcomes}
Dupa capitolul asta vei \textcommabelow{s}ti:
\begin{itemize}[itemsep=1pt]
  \item Configureaza \code{<title>} \textcommabelow{s}i \code{<meta description>} pentru rezultatele de cautare.
  \item Adauga cele 5 tag-uri \code{og:*} pentru previzualizare in re\textcommabelow{t}ele sociale.
  \item Genereaza \textcommabelow{s}i valideaza imaginea \code{og:image} la dimensiunile \textcommabelow{s}i formatul corect.
  \item Adauga date structurate (JSON-LD) pentru rich snippets in Google.
  \item Genereaza \code{robots.txt} \textcommabelow{s}i \code{sitemap.xml} pentru a ghida crawler-ele.
  \item Distinge mit-urile SEO de practica actuala.
  \item \faRocket\ \emph{Bonus self-study:} tipuri Schema.org (Article, BreadcrumbList, FAQ), \code{preconnect}/\code{dns-prefetch}, linking intern, mobile-first indexing.
\end{itemize}
\end{outcomes}

\begin{whymatters}
Un site fara SEO este invizibil. Aplica\textcommabelow{t}iile sociale moderne (WhatsApp, Discord, LinkedIn) afi\textcommabelow{s}eaza by default cardul de previzualizare; un site fara meta tags arata neprofesional cand cineva il share-uieste. SEO de baza este \emph{configurare}, nu \emph{magic}: cele 7\textendash{}10 tag-uri prezentate aici acopera 80\% din optimizarea esentiala.
\end{whymatters}

\begin{nota}[Pre-rechizite]
Capitolul presupune un site deployat (Faza 4) cu URL public accesibil. Tag-urile se valideaza pe varianta live, nu pe localhost.
\end{nota}

% =========================================================
\section{Doua audiente, doua seturi de tag-uri}

\begin{itemize}
  \item \textbf{Motoare de cautare} (Google, Bing). Folosesc \code{<title>}, \code{<meta name="description">}, structura semantica, datele Schema.org.
  \item \textbf{Platforme sociale} (Facebook, X, LinkedIn, WhatsApp, Discord). Folosesc protocolul Open Graph: \code{<meta property="og:*">}.
\end{itemize}

% =========================================================
\section{Metadate pentru motoarele de cautare}

Doua elemente determina aspectul paginii in rezultatele cautarii:

\begin{lstlisting}[style=html, caption={Plasare in <head>}]
<title>Voluntar BEST -- Student CS, Iasi</title>
<meta name="description" content="Pagina personala a Voluntarului BEST, studenta la CS, pasionata de web development.">
\end{lstlisting}

\subsection{Calitatea \texttt{<title>}}

\begin{dano}[Asa nu]
\begin{lstlisting}[style=html, numbers=none]
<title>Document</title>
<title>Untitled</title>
<title>Index</title>
<title>Acasa</title>
<title>web web web html portofoliu cv student iasi cs uaic</title>
\end{lstlisting}
\textbf{Probleme.} Generic, nesemnificativ; sau spam de cuvinte-cheie penalizat de Google.
\end{dano}

\begin{dado}[Asa da]
\begin{lstlisting}[style=html, numbers=none]
<title>Voluntar BEST -- Portofoliu Web Developer, Iasi</title>
<title>Despre BEST Iasi -- Asociatie studenteasca tehnica</title>
\end{lstlisting}
\textbf{Reguli.} 50\textendash{}60 caractere; specific paginii (nu site-ului); structura ,,Subiect \textendash{} Context''.
\end{dado}

\subsection{Calitatea \texttt{<meta description>}}

Aceleasi reguli ca pentru \code{<title>}, doar lungimea difera: \textbf{150\textendash{}160 caractere}, descrie con\textcommabelow{t}inutul concret (proiecte, CV, sec\textcommabelow{t}iuni), fara fraze de umplutura tip \emph{,,Bine a\textcommabelow{t}i venit''}.

\begin{lstlisting}[style=html, numbers=none]
<meta name="description" content="Portofoliul Voluntarului BEST: proiecte web (HTML, CSS, JS), CV, contact. Voluntar BEST Iasi si student CS la UAIC.">
\end{lstlisting}

% =========================================================
\section{Protocolul Open Graph}

Open Graph Protocol (OGP), specificat de Facebook in 2010, este standardul de fapt pentru previzualizarea paginilor web in retelele sociale. Implementarea consta in adaugarea a cinci tag-uri \code{<meta>} cu prefixul \code{og:}:

\begin{lstlisting}[style=html, caption={Tag-uri OGP minimale}]
<meta property="og:type" content="website">
<meta property="og:title" content="Voluntar BEST -- Student CS, Iasi">
<meta property="og:description" content="Pagina personala. Web dev, design, open source.">
<meta property="og:image" content="https://voluntar-best.github.io/site/og.png">
<meta property="og:url" content="https://voluntar-best.github.io/site/">
\end{lstlisting}

\begin{itemize}
  \item \code{og:type} \textendash{} categoria continutului. Valoarea \code{website} acopera majoritatea cazurilor; alternative: \code{article}, \code{video}, \code{book}, \code{profile}.
  \item \code{og:title} si \code{og:description} \textendash{} pot diferi de \code{<title>} si \code{<meta description>}, fiind optimizate specific pentru context social (mai punctuale, cu CTA).
  \item \code{og:image} \textendash{} dimensiunea recomandata 1200\,x\,630 pixeli (raport 1.91:1).
  \item \code{og:url} \textendash{} URL-ul canonic al paginii.
\end{itemize}

\subsection{Constrangerea \texttt{og:image}: URL absolut}

\begin{dano}[Asa nu \textendash{} URL relativ]
\code{<meta property="og:image" content="og.png">}\\
\code{<meta property="og:image" content="/images/og.png">}\\
\textbf{Problema.} Crawler-ele platformelor sociale nu pot rezolva caile relative, intrucat nu cunosc contextul de baza al paginii sursa. Cardul afisat ramane fara imagine.
\end{dano}

\begin{dado}[Asa da \textendash{} URL absolut]
\code{<meta property="og:image" content="https://voluntar-best.github.io/site/og.png">}\\
\textbf{Beneficiu.} Imaginea este accesibila direct prin URL si poate fi cache-uita de Facebook, Slack, Discord.
\end{dado}

\subsection{Calitatea imaginii \texttt{og:image}}

\begin{itemize}[itemsep=2pt]
  \item Dimensiune ideala: \textbf{1200\,x\,630 px} (raport 1.91:1, formatul afisat de Facebook si LinkedIn).
  \item Dimensiune minima: 600\,x\,315 px. Sub aceasta limita, unele platforme nu afiseaza cardul mare.
  \item Format: JPG sau PNG; dimensiunea maxima 8 MB (Facebook), in practica $\leq$1 MB pentru incarcare rapida.
  \item Continut lizibil la dimensiune redusa: titlul + 1\textendash{}2 elemente vizuale, nu paragrafe intregi.
\end{itemize}

% =========================================================
\section{Twitter Card}

Platforma X (anterior Twitter) suporta nativ OGP, dar accepta si tag-uri proprii pentru carduri cu imagine mare:

\begin{lstlisting}[style=html]
<meta name="twitter:card" content="summary_large_image">
\end{lstlisting}

Adaugarea acestei singure linii instruieste X sa afiseze cardul in formatul extins, folosind valorile deja prezente in tag-urile \code{og:*}.

% =========================================================
\section{Date structurate (JSON-LD)}

Schema.org defineste un vocabular standardizat pentru descrierea entitatilor (persoane, organizatii, evenimente, articole) intr-un format pe care motoarele de cautare il proceseaza explicit. Implementarea recomandata: JSON-LD plasat intr-un \code{<script>} cu tipul \code{application/ld+json}.

\begin{lstlisting}[style=html, caption={Schema Person pentru o pagina personala}]
<script type="application/ld+json">
{
  "@context": "https://schema.org",
  "@type": "Person",
  "name": "Voluntar BEST",
  "url": "https://voluntar-best.github.io/site/",
  "jobTitle": "Student CS",
  "alumniOf": "Universitatea Alexandru Ioan Cuza",
  "sameAs": [
    "https://github.com/voluntar-best",
    "https://www.linkedin.com/in/voluntar-best"
  ]
}
</script>
\end{lstlisting}

Avantajul principal: Google poate afisa pagina cu \emph{rich snippets} (foto, link-uri sociale, descriere structurata) in rezultatele cautarii. Validarea se face cu Google Rich Results Test (\url{https://search.google.com/test/rich-results}).

\begin{ponturi}
Pentru un eveniment, foloseste \code{"@type": "Event"} cu campurile \code{startDate}, \code{location}, \code{organizer}. Google poate afisa evenimentul direct intr-un panel lateral. Vocabularul complet: \url{https://schema.org}.
\end{ponturi}

% =========================================================
\section{URL-ul canonic}

Cand acelasi continut este accesibil prin mai multe URL-uri (cu si fara \code{www}, cu si fara slash final, cu parametri UTM, etc.), Google trebuie sa stie care URL este versiunea ,,oficiala''. Tag-ul canonical previne penalizarea pentru continut duplicat:

\begin{lstlisting}[style=html]
<link rel="canonical" href="https://voluntar-best.github.io/site/">
\end{lstlisting}

Indicatie: pe paginile de proiect, URL-ul canonic este URL-ul ,,curat'' al paginii, fara parametri de tracking.

% =========================================================
\section{Mituri SEO}

\begin{itemize}
  \item \textbf{Mit:} \code{<meta name="keywords">} influenteaza ranking-ul.
  \item \textbf{Realitate:} Google a anuntat oficial in 2009 ca ignora complet acest tag. Bing si Yandex il folosesc minim. \textbf{A nu se mai folosi.}
\end{itemize}

\begin{itemize}
  \item \textbf{Mit:} repetarea cuvintelor cheie in continut creste ranking-ul.
  \item \textbf{Realitate:} \emph{keyword stuffing} este penalizat. Algoritmii moderni (BERT, MUM) inteleg sensul textului si premiaza scrierea naturala.
\end{itemize}

\begin{itemize}
  \item \textbf{Mit:} site-urile noi sunt penalizate de Google.
  \item \textbf{Realitate:} indexarea dureaza 1\textendash{}7 zile pentru un site nou. Inscrierea in Google Search Console (\url{https://search.google.com/search-console}) accelereaza procesul.
\end{itemize}

% =========================================================
\section{\texttt{robots.txt} \textendash{} ghideaza crawler-ele}

\code{robots.txt} este un fisier text plasat in radacina site-ului (\code{https://voluntar-best.github.io/robots.txt}) care indica robotilor de crawl ce pagini sa indexeze si unde sa caute sitemap-ul. Format strict, dar simplu:

\begin{lstlisting}[caption={\texttt{robots.txt} minimal pentru GitHub Pages}, basicstyle=\ttfamily\small, backgroundcolor=\color{codeBg}, frame=leftline, framerule=2.5pt, rulecolor=\color{partcolor}, framexleftmargin=4pt, xleftmargin=14pt]
User-agent: *
Allow: /

Sitemap: https://voluntar-best.github.io/sitemap.xml
\end{lstlisting}

Cum se cite\textcommabelow{s}te:
\begin{itemize}[itemsep=1pt]
  \item \code{User-agent: *} \textendash{} regulile se aplica oricarui robot (Googlebot, Bingbot, etc.).
  \item \code{Allow: /} \textendash{} permite accesul la intreg site-ul.
  \item \code{Disallow: /admin/} (alternativ) \textendash{} interzice anumite cai. Util pentru pagini in lucru, dashboard-uri private.
  \item \code{Sitemap: ...} \textendash{} URL absolut catre sitemap-ul XML (vezi sectiunea urmatoare).
\end{itemize}

\begin{atentie}
\code{robots.txt} este o \emph{conventie}, nu un mecanism de securitate. Crawler-ele rau-voitoare il ignora. Daca ai date sensibile, foloseste autentificare reala, nu \code{Disallow}.
\end{atentie}

\begin{ponturi}
Verificare: deschide direct in browser \code{https://<site>/robots.txt}. Daca se afiseaza textul, e plasat corect.
\end{ponturi}

% =========================================================
\section{\texttt{sitemap.xml} \textendash{} ce pagini exista}

Sitemap-ul este un fisier XML care listeaza toate paginile importante ale site-ului. Google il foloseste pentru indexare mai rapida si completa, in special pentru site-urile noi sau cu structura adanca.

\begin{lstlisting}[caption={\texttt{sitemap.xml} scris manual pentru un portofoliu mic}, basicstyle=\ttfamily\small, backgroundcolor=\color{codeBg}, frame=leftline, framerule=2.5pt, rulecolor=\color{partcolor}, framexleftmargin=4pt, xleftmargin=14pt]
<?xml version="1.0" encoding="UTF-8"?>
<urlset xmlns="http://www.sitemaps.org/schemas/sitemap/0.9">
  <url>
    <loc>https://voluntar-best.github.io/</loc>
    <lastmod>2026-05-11</lastmod>
    <priority>1.0</priority>
  </url>
  <url>
    <loc>https://voluntar-best.github.io/projects.html</loc>
    <lastmod>2026-05-08</lastmod>
    <priority>0.8</priority>
  </url>
  <url>
    <loc>https://voluntar-best.github.io/contact.html</loc>
    <lastmod>2026-04-22</lastmod>
    <priority>0.5</priority>
  </url>
</urlset>
\end{lstlisting}

Reguli:
\begin{itemize}[itemsep=1pt]
  \item \code{<loc>} \textendash{} URL absolut (obligatoriu).
  \item \code{<lastmod>} \textendash{} data ultimei modificari, format ISO \code{YYYY-MM-DD} (op\textcommabelow{t}ional, recomandat).
  \item \code{<priority>} \textendash{} 0.0 \textendash{} 1.0, relativ in cadrul site-ului. Home = 1.0, pagini interne = 0.5\textendash{}0.8.
  \item Pentru site-uri mici (\textless{}50 pagini), scrierea manuala este OK. Pentru site-uri mari, foloseste generatoare automate.
\end{itemize}

\subsection{Submisie in Google Search Console}

\begin{enumerate}[itemsep=1pt]
  \item Deschide \url{https://search.google.com/search-console}, adauga proprietatea (URL prefix: \code{https://voluntar-best.github.io/}).
  \item Verifica posesiunea \textendash{} pentru GitHub Pages, metoda recomandata e meta tag in \code{<head>}.
  \item Sec\textcommabelow{t}iunea ,,Sitemaps'' $\to$ adauga \code{sitemap.xml} $\to$ Submit.
  \item In 1\textendash{}7 zile, Google va incepe sa indexeze toate URL-urile listate.
\end{enumerate}

\begin{ponturi}
\emph{Beneficiul real:} fara sitemap, Google poate descoperi pagini doar prin link-uri. Cu sitemap, le \textbf{declari explicit}. Pentru un site nou (fara backlink-uri inca), diferen\textcommabelow{t}a e intre indexare in 7 zile vs in 1\textendash{}2 zile.
\end{ponturi}

% =========================================================
\section{Favicon}

Favicon-ul (iconita afisata in tab-ul browserului si in bookmark-uri) se declara prin elementul \code{<link>}:

\begin{lstlisting}[style=html]
<link rel="icon" href="favicon.svg" type="image/svg+xml">
<link rel="apple-touch-icon" href="apple-touch-icon.png">
\end{lstlisting}

Generarea unui favicon simplu se poate face cu serviciul \url{https://favicon.io/favicon-generator/}, care produce fisierele necesare pornind de la o singura litera si o culoare. Pentru un favicon profesional din SVG, conversia automata in toate dimensiunile necesare se face cu \url{https://realfavicongenerator.net}.

% =========================================================
\begin{bonusbox}
\textbf{Acest box e optional.} Daca esti la inceput, treci direct la \emph{Verificare}. Mai jos sunt pattern-urile pe care le folosesc site-urile cu trafic real.

\medskip

\textbf{\faSitemap\ Tipuri Schema.org dincolo de Person}

Vocabularul Schema.org are sute de \code{@type}-uri. Cele mai utile in practica:

\medskip

\emph{Article \textendash{} pentru blog/posts:}
\begin{lstlisting}[style=html, numbers=none]
{
  "@context": "https://schema.org",
  "@type": "Article",
  "headline": "Cum am invatat web development in 2 saptamani",
  "author": { "@type": "Person", "name": "Voluntar BEST" },
  "datePublished": "2026-05-11",
  "dateModified": "2026-05-11",
  "image": "https://voluntar-best.github.io/blog/learning/cover.jpg"
}
\end{lstlisting}

\emph{BreadcrumbList \textendash{} navigare ierarhic\u a in rezultate:}
\begin{lstlisting}[style=html, numbers=none]
{
  "@context": "https://schema.org",
  "@type": "BreadcrumbList",
  "itemListElement": [
    { "@type": "ListItem", "position": 1, "name": "Acasa",     "item": "https://voluntar-best.github.io/" },
    { "@type": "ListItem", "position": 2, "name": "Proiecte",  "item": "https://voluntar-best.github.io/projects/" },
    { "@type": "ListItem", "position": 3, "name": "Portfolio Web" }
  ]
}
\end{lstlisting}
Google afi\textcommabelow{s}eaza aceasta ierarhie direct in rezultatele de cautare \textendash{} click rate semnificativ mai bun.

\emph{FAQPage \textendash{} intreb\u ari frecvente, afi\textcommabelow{s}ate ca rich snippet:}
\begin{lstlisting}[style=html, numbers=none]
{
  "@context": "https://schema.org",
  "@type": "FAQPage",
  "mainEntity": [{
    "@type": "Question",
    "name": "Cum aplic la BEST Iasi?",
    "acceptedAnswer": { "@type": "Answer", "text": "Inscrie-te la recrutarea din toamna pe best.org.ro." }
  }]
}
\end{lstlisting}

\emph{Event \textendash{} cu datele cheie indexate:}
\begin{lstlisting}[style=html, numbers=none]
{
  "@context": "https://schema.org",
  "@type": "Event",
  "name": "BEST Training Web 2026",
  "startDate": "2026-11-15T10:00:00+02:00",
  "location": { "@type": "Place", "name": "TUIASI" },
  "organizer": { "@type": "Organization", "name": "BEST Iasi" }
}
\end{lstlisting}

\medskip

\textbf{\faPlug\ \texttt{preconnect} \textcommabelow{s}i \texttt{dns-prefetch} \textendash{} performance + SEO signal}

Cand pagina ta foloseste resurse de pe alt domeniu (Google Fonts, CDN, Analytics), DNS lookup + TCP handshake adauga 100\textendash{}300ms inainte ca resursa sa inceapa sa se descarce. \code{preconnect} face acest setup in paralel cu HTML parsing, in fundal:

\begin{lstlisting}[style=html, numbers=none]
<link rel="preconnect" href="https://fonts.googleapis.com">
<link rel="preconnect" href="https://fonts.gstatic.com" crossorigin>
<link rel="dns-prefetch" href="https://www.google-analytics.com">
\end{lstlisting}

\code{preconnect} \textendash{} DNS + TCP + TLS (cost mai mare, dar resursa e gata sa fie descarcata).\\
\code{dns-prefetch} \textendash{} doar DNS (mai ieftin, foloseste pentru \textgreater 3 domenii).

Impact: imbunata\textcommabelow{t}e\textcommabelow{s}te LCP (Core Web Vitals) cu 100\textendash{}500ms cand pagina depinde de un CDN.

\medskip

\textbf{\faLink\ Linking intern \textendash{} arhitectur\u a, nu trick}

Google urmare\textcommabelow{s}te link-urile interne pentru a in\textcommabelow{t}elege relevan\textcommabelow{t}a relativa a paginilor. Pagini cu mai multe link-uri intrand sunt considerate ,,mai importante''.

\begin{itemize}[itemsep=1pt]
  \item Folose\textcommabelow{s}te \textbf{anchor text descriptiv} (\code{<a href="/projects">proiectele mele web</a>}), nu \code{click aici}.
  \item Conecteaza paginile inrudite: pe pagina de proiect, link spre alte proiecte. Pe blog post, link catre articole inrudite.
  \item Evita \textbf{orphan pages} (pagini fara niciun link intrand). Google le indexeaza tarziu sau deloc.
\end{itemize}

\medskip

\textbf{\faMobile*\ Mobile-first indexing}

Din 2023, Google indexeaza \textbf{varianta mobila} a paginii, nu desktop. Pagini care arata bine pe desktop dar sunt rupte pe mobile sunt penalizate masiv. Verifica:

\begin{itemize}[itemsep=1pt]
  \item \code{<meta name="viewport" content="width=device-width, initial-scale=1">} \textendash{} obligatoriu.
  \item Textul lizibil fara zoom (\code{font-size} \textgreater{} 14px pe mobil).
  \item Butoane clic\u abile (\textgreater 44x44px, conform Apple HIG).
  \item Con\textcommabelow{t}inutul important \textbf{nu este ascuns} pe mobil (Google considera continutul vizibil ca cel canonic).
  \item Aceea\textcommabelow{s}i pagina pe mobil \textcommabelow{s}i desktop \textendash{} nu varianta separata \code{m.site.ro}.
\end{itemize}

\medskip

\textbf{\faChartBar\ Performance ca semnal de ranking}

Din 2021, Core Web Vitals (LCP, INP, CLS \textendash{} vezi Faza 6) sunt \textbf{semnal de ranking}. Doua pagini cu con\textcommabelow{t}inut similar: cea mai rapida castiga. Nu vei concura cu site-uri lente.

Recap rapid:
\begin{itemize}[itemsep=1pt]
  \item \textbf{LCP} \textless 2.5s \textendash{} cea mai mare imagine/text apare repede.
  \item \textbf{INP} \textless 200ms \textendash{} pagina raspunde rapid la click/tap.
  \item \textbf{CLS} \textless 0.1 \textendash{} layout-ul nu ,,sare'' dupa incarcare.
\end{itemize}

\medskip

\textbf{\faRoute\ Unde mergi de aici}

\begin{itemize}[itemsep=1pt]
  \item \emph{Schema.org vocabular complet:} \url{https://schema.org/docs/full.html}
  \item \emph{Generator JSON-LD vizual:} \url{https://technicalseo.com/tools/schema-markup-generator/}
  \item \emph{Google Mobile-Friendly Test:} \url{https://search.google.com/test/mobile-friendly}
  \item \emph{Sitemap generator pentru site-uri statice:} \url{https://www.xml-sitemaps.com}
\end{itemize}
\end{bonusbox}

% =========================================================
\section*{Verificare cuno\textcommabelow{s}tin\textcommabelow{t}e}

\begin{selfcheck}
\begin{enumerate}[itemsep=2pt]
  \item Cate tag-uri \code{og:*} sunt minim necesare pentru un card valid de previzualizare?
  \item De ce \code{og:image} trebuie sa fie URL absolut, nu relativ?
  \item \emph{Spot the issue:} \code{<title>Document</title>}.
  \item Care este diferen\textcommabelow{t}a func\textcommabelow{t}ionala intre \code{<title>} si \code{og:title}?
  \item De ce \code{<meta name="keywords">} este considerat depasit?
  \item Ce reprezinta directiva \code{Sitemap:} in \code{robots.txt}?
  \item De ce \code{<lastmod>} in \code{sitemap.xml} este util pentru Google?
\end{enumerate}
\end{selfcheck}

% =========================================================
\section*{Recapitulare}

\begin{recap}[Faza 5 -- SEO + Open Graph]
\begin{itemize}[itemsep=1pt]
  \item Doua audiente: motoare de cautare (\code{<title>} + \code{<meta description>}) si retele sociale (5 \code{og:*} tags).
  \item \code{<title>}: 50\textendash{}60 caractere, specific, structura ,,Subiect \textendash{} Context''.
  \item \code{<meta description>}: 150\textendash{}160 caractere, descriere concreta, fara umplutura.
  \item \code{og:image}: URL absolut, dimensiune optima 1200\,x\,630, format JPG/PNG, $\leq$1MB.
  \item \code{<meta name="twitter:card" content="summary\_large\_image">} activeaza cardul mare pe X.
  \item JSON-LD (Schema.org) descrie entitati: Person, Organization, Article, Event.
  \item \code{<link rel="canonical">} previne penalizarea pentru continut duplicat.
  \item \code{robots.txt} ghideaza crawler-ele \textcommabelow{s}i declara sitemap-ul. \code{sitemap.xml} listeaza URL-urile importante.
  \item Submisia in Google Search Console accelereaza indexarea pentru site-uri noi.
  \item \code{<meta name="keywords">} este ignorat de Google din 2009 \textendash{} \textbf{a nu se folosi}.
  \item \faRocket\ Bonus: tipuri Schema.org (Article, BreadcrumbList, FAQ, Event), \code{preconnect}/\code{dns-prefetch}, linking intern, mobile-first indexing.
\end{itemize}
\end{recap}

\partdivider

% =========================================================
\begin{joc}
\textbf{Platforma:} OpenGraph.xyz (\url{https://opengraph.xyz}).

\medskip

\textbf{Activitatea.}
\begin{enumerate}[itemsep=1pt]
  \item Adaugarea celor cinci tag-uri \code{og:*} in \code{<head>}, cu valori specifice paginii personale.
  \item Inlocuirea valorii pentru \code{og:image} cu un URL absolut catre o imagine 1200\,x\,630, accesibila public.
  \item Commit si push: \code{git add .; git commit -m "Adauga meta tags Open Graph"; git push}.
  \item Accesarea OpenGraph.xyz si lipirea URL-ului public al paginii.
  \item Verificarea previzualizarilor pentru Facebook, X, LinkedIn, WhatsApp, iMessage.
\end{enumerate}

\textbf{Diagnoza problemelor:} (1) URL-ul \code{og:image} nu este absolut; (2) imaginea returneaza HTTP 404; (3) dimensiunea imaginii este sub 600 pixeli pe oricare axa.

\medskip

\textbf{Durata estimata:} 5 minute.
\end{joc}

% =========================================================
\section*{\faGraduationCap\ \ Resurse pentru aprofundare}

\begin{itemize}[itemsep=2pt]
  \item Open Graph Protocol. \emph{Specificatia oficiala}. \url{https://ogp.me}
  \item Meta Tags. \emph{Generator vizual}. \url{https://metatags.io}
  \item Google. \emph{Rich Results Test}. \url{https://search.google.com/test/rich-results}
  \item Google. \emph{Search Console} (indexare manuala). \url{https://search.google.com/search-console}
  \item Google. \emph{Document metadata}. \url{https://web.dev/learn/html/metadata}
  \item Schema.org. \emph{Structured data documentation}. \url{https://schema.org/docs/gs.html}
  \item X Inc. \emph{Twitter Cards documentation}. \url{https://developer.x.com/en/docs/twitter-for-websites/cards}
\end{itemize}

\bigskip

\bridge{Capitolul 6 (Lighthouse) ofera o masuratoare obiectiva a calita\textcommabelow{t}ii. Toate fix-urile aplicate in fazele 1\textendash{}5 sunt acum cuantificate intr-un scor 0\textendash{}100.}
